% ============================================================================
% Fichier  : front/02_mode_emploi.tex — Mode d'emploi
% Version  : 1.0 -- novembre 2026
% ============================================================================
% Règle des diapositives (GCI, ACSI) : 30 h en présentiel = 15 séances de
% 2 h ; fiches et ateliers hors des 30 h. Répartition des séances indicative.
% ============================================================================

\chapter*{Mode d'emploi du cours}
\addcontentsline{toc}{chapter}{Mode d'emploi du cours}
\markboth{Mode d'emploi}{Mode d'emploi}

\section*{Structure générale}

Le cours compte un prologue et six parties, soit 16~chapitres, pour
30~heures en présentiel réparties en quinze séances de deux heures, qui
couvrent l'ensemble des parties. Les fiches de questions-réponses par
chapitre et les dix laboratoires de 60~minutes de NDAP SUKULU sont des
travaux à faire à la maison, hors de ces 30~heures. Les corrigés des
exercices sont réunis dans un livret séparé.

\begin{tableaucours}{cL{6.4cm}cL{3.2cm}}{%
    \tableauHeader{Partie} &
    \tableauHeader{Thème} &
    \tableauHeader{Chap.} &
    \tableauHeader{Niveau}
}
0 & Prologue~: certifiée, et pourtant compromise & P & introductif \\
I & Fondements~: gouvernance, risque et incertitude & 1--2 & fondamental \\
II & Méthodes~: gérer les risques, conduire l'audit & 3--7 & fondamental à
intermédiaire \\
III & Cadre normatif et juridique & 8--10 & fondamental \\
IV & Pratique~: SMSI, audit technique et GRC & 11--14 & intermédiaire à
expert \\
V & Défis avancés & 15 & expert \\
VI & Recherche~: mesurer ce que l'on ne sait pas & 16 & recherche \\
\end{tableaucours}

\section*{Les quinze séances}

\begin{tableaucours}{cL{6.8cm}L{4.2cm}}{%
    \tableauHeader{Séance} &
    \tableauHeader{Contenu} &
    \tableauHeader{Travail à la maison}
}
1 & Prologue~: le certificat, le périmètre, l'inconnu & Fiche P \\
2 & Chap.~1~: gouverner la cybersécurité & Fiche 1 \\
3 & Chap.~2~: le risque cyber et sa cote & Fiche 2, atelier TD01 \\
4 & Chap.~3~: ISO~31000 et ISO/IEC~27005 & Fiche 3 \\
5 & Chap.~4~: EBIOS Risk Manager & Fiche 4, ateliers TD02 et TD03 \\
6 & Chap.~5~: quantifier avec FAIR & Fiche 5, atelier TD04 \\
7 & Chap.~6~: conduire un audit & Fiche 6 \\
8 & Chap.~7~: décomposer, coter, conclure & Fiche 7, atelier TD05 \\
9 & Chap.~8~: la famille ISO/IEC~27000 & Fiche 8, atelier TD06 \\
10 & Chap.~9~: référentiels et correspondances & Fiche 9, atelier TD07 \\
11 & Chap.~10~: droit et déontologie & Fiche 10 \\
12 & Chap.~11~: mettre en œuvre et auditer un SMSI & Fiche 11, atelier TD08 \\
13 & Chap.~12~: l'audit technique & Fiche 12, atelier TD09 \\
14 & Chap.~13 et~14~: la GRC en pratique, les cas intégrateurs & Fiches 13
et 14, atelier TD10 \\
15 & Chap.~15 et~16~: risques émergents, mesurer l'inconnu & Fiches 15 et
16 \\
\end{tableaucours}

Cette répartition est indicative~; l'enseignant l'adapte au calendrier et au
rythme du groupe.

\section*{Les dix laboratoires de NDAP SUKULU}

Les dix laboratoires, NRAC-LAB-01 à NRAC-LAB-10, sont présentés dans la
section « NDAP SUKULU, la salle de classe compagnon », qui suit ce mode
d'emploi.

\section*{Le gabarit des chapitres}

Chaque chapitre suit le même plan~: une épigraphe, le niveau et les
prérequis, les objectifs, une situation d'ouverture tirée de l'un des trois
cas, cinq à huit sections, ce qu'il faut retenir, des points de contrôle,
trois à quatre exercices gradués, des lectures, et une transition.

\section*{Les encadrés}

\begin{tableaucours}{L{3.6cm}L{8.6cm}}{%
    \tableauHeader{Encadré} &
    \tableauHeader{Rôle}
}
Objectifs & Ce que l'on doit savoir faire à la fin du chapitre \\
Définition & Les notions, en gras dans le texte et indexées \\
Violation typique & Une erreur fréquente, et la règle qui l'évite \\
Information & Un complément, un renvoi, une nuance \\
Opération (SANAGA, LIZBIZ, ANPP) & Le cas~: les constatations, puis
l'analyse, avec un niveau de confiance \\
Lecture & Un calcul de $\mu$, $U$ et $H$, et sa lecture \\
Fiche de risque & Un risque décrit, coté et décidé \\
Papier de travail & Ce que le dossier doit contenir pour être refait par un
autre \\
Métier & Ce que la notion change pour chacun des cinq métiers du cours \\
Texte juridique & Un article cité dans son texte officiel, et sa portée \\
Ce qu'il faut retenir & La synthèse du chapitre \\
Exercice & Trois niveaux de difficulté~: \niveaubase, \niveaumoyen,
\niveauexpert \\
\end{tableaucours}

\section*{Les conventions}

\begin{itemize}
    \item \textbf{Les cotes}~: \OUI{} (établi par une preuve), \NON{}
          (absent ou contredit par une preuve), \NEK{} (connaissance
          insuffisante). Aucun barème ne donne de demi-point à \NEK{}.
    \item \textbf{Les mesures}~: le taux de conformité $\mu$, le taux
          d'incertitude $U$ et l'entropie $H$, en bits, se lisent toujours
          ensemble~; règle de décision~: $U$ d'abord, $\mu$ ensuite, $H$ pour
          contrôler (Chapitre~\ref{chap:conclure}).
    \item \textbf{L'échelle de cotation} des risques est celle du
          Chapitre~\ref{chap:risque} (quatre niveaux de vraisemblance et
          d'impact, criticité de 1 à 16).
    \item \textbf{Constatation et analyse} sont toujours séparées~; l'analyse
          porte un niveau de confiance.
    \item \textbf{Les textes}~: seuls les articles lus dans leur version
          officielle sont cités entre guillemets~; les normes sous licence
          sont paraphrasées~; le degré de vérification figure à
          l'annexe~\ref{ann:textes}.
    \item \textbf{Les travaux de recherche} sont cités avec leur statut~:
          prépublication déposée, ou travail en préparation
          (Chapitre~\ref{chap:auditologie}).
    \item \textbf{Les cas} sont fictifs~: noms en majuscules, montants en
          francs CFA, villes réelles.
\end{itemize}
